% ⚠️ UNAPPROVED: every statement in this file was drafted by Claude (Fable 5.2) on 2026-11-10 and has NOT been reviewed by Nathan.
%
% AI-derived (Claude, 2026-11-10), unverified by Nathan.
% Draft replacement for the passage around eq:hsi-average-error-bound and eq:hsi-convergence
% in "ALT format.tex" (Section "Optimality properties"). Uses the paper's macros (\Eta, \coloneq).
% Lean counterpart: HsiLean/Rate.lean (all statements there are UNAPPROVED too). The Lean name
% of each step is given in a comment. The Lean proofs check with no sorry and only the standard
% axioms (propext, Classical.choice, Quot.sound).
%
% WHAT CHANGES RELATIVE TO THE PAPER (please check each):
%  (a) Two hypotheses are made explicit (Assumption below). Without (ii) the arrow in
%      eq:hsi-convergence is FALSE as stated: see the counterexample remark at the end.
%  (b) eq:hsi-error-bound is stated with \le, not <. (Strictness needs a second model with
%      positive weight and positive probability on X; it plays no role in the rate.)
%  (c) The expectation in eq:hsi-convergence is taken over the whole dataset X as well as the
%      new sequence x_nu, not conditionally on a fixed X.
%  (d) The hyperprior weights may sum to less than 1; the chain rule then has the extra
%      constant +log_2 sum_j Eta(U_j) <= 0, which only helps.

\paragraph{Setting.}
Let $A\subseteq\mathbb{B}^*$ be the set of sequences that may appear in a dataset, and let the ordered dataset $X=(x_1,x_2,\dots)$ be sampled independently from $\mu$, with $x_{<\nu}\coloneq(x_1,\dots,x_{\nu-1})$.

\begin{assumption}
\label{ass:hsi-rate}
(i) $\mu$ is a probability distribution on $A$, that is, $\sum_{x\in A}P_{U_\mu}(x)=1$, and $\Eta(U_\mu)>0$.
(ii) $\sum_{x\in A}P_{U_j}(x)\le 1$ for every $j$.
\end{assumption}
Condition (ii) holds for every semimeasure whenever $A$ is prefix-free, for example when all sequences in the dataset have the same length; it can fail if $A$ contains a sequence together with one of its extensions.

Write the conditioned predictor in normalised form,
\[
    P_{\Eta(\cdot\mid X)}(x) \coloneq \frac{P_\Eta(X,x)}{P_\Eta(X)}
    = \frac{\sum_j \Eta(U_j\mid X)\,P_{U_j}(x)}{\sum_j \Eta(U_j\mid X)},
\]
and its excess error on the $\nu$-th sequence as
\[
    e_\nu \coloneq -\log_2\frac{P_{\Eta(\cdot\mid x_{<\nu})}(x_\nu)}{P_{U_\mu}(x_\nu)} .
\]

\begin{lemma}[Chain rule over sequences]
\label{lem:hsi-chain-rule}
% Lean: sum_excess_eq_totalExcess
For every dataset $X$ with $P_{U_\mu}(X)>0$,
\begin{equation}
\label{eq:hsi-chain-rule}
    \sum_{\nu=1}^{|X|} e_\nu
    = -\log_2\frac{P_\Eta(X)}{P_{U_\mu}(X)} + \log_2\sum_j\Eta(U_j).
\end{equation}
The last term is $0$ if the hyperprior sums to one and is otherwise negative.
\end{lemma}
\begin{proof}
Since $P_{U_\mu}(X)=\prod_\nu P_{U_\mu}(x_\nu)$ and $P_\Eta(X)\ge\Eta(U_\mu)P_{U_\mu}(X)>0$, the product
\[
    \prod_{\nu=1}^{|X|}\frac{P_{\Eta(\cdot\mid x_{<\nu})}(x_\nu)}{P_{U_\mu}(x_\nu)}
    =\frac{1}{P_{U_\mu}(X)}\prod_{\nu=1}^{|X|}\frac{P_\Eta(x_{\le\nu})}{P_\Eta(x_{<\nu})}
    =\frac{P_\Eta(X)}{P_{U_\mu}(X)\,P_\Eta(\emptyset)}
\]
telescopes, with $P_\Eta(\emptyset)=\sum_j\Eta(U_j)$. Take $-\log_2$.
\end{proof}

\begin{theorem}[Rate of convergence]
\label{thm:hsi-rate}
Under Assumption \ref{ass:hsi-rate}:
\begin{enumerate}
    \item % Lean: sum_excess_le
    $\sum_{\nu=1}^{n} e_\nu \le -\log_2\Eta(U_\mu)$ for every $n$, with $\mu$-probability one.
    \item % Lean: integrable_excessAt, expected_excess_nonneg
    $E_\mu(e_\nu)$ is finite and $E_\mu(e_\nu)\ge 0$ for every $\nu$.
    \item % Lean: sum_expected_excess_le, tsum_expected_excess_le
    $\sum_{\nu=1}^{\infty} E_\mu(e_\nu)\le -\log_2\Eta(U_\mu)$.
    \item % Lean: expected_excess_rate
    For every $\varepsilon>0$,
    \begin{equation}
    \label{eq:hsi-rate}
        \bigl|\{\nu : E_\mu(e_\nu)>\varepsilon\}\bigr| \;\le\; \frac{-\log_2\Eta(U_\mu)}{\varepsilon},
    \end{equation}
    independently of $|X|$.
    \item % Lean: tendsto_expected_excess
    In particular $E_\mu(e_\nu)\to 0$ as $\nu\to\infty$, which is Equation \ref{eq:hsi-convergence}.
\end{enumerate}
\end{theorem}
\begin{proof}
(1) With probability one every sampled $x_\nu$ has $P_{U_\mu}(x_\nu)>0$, so Lemma \ref{lem:hsi-chain-rule} applies to $x_{\le n}$, and Equation \ref{eq:hsi-error-bound} bounds its right-hand side by $-\log_2\Eta(U_\mu)$.

(2) Fix $x_{<\nu}$ and write $q(x)\coloneq P_{\Eta(\cdot\mid x_{<\nu})}(x)$. By Assumption \ref{ass:hsi-rate}(ii),
\[
    \sum_{x\in A}q(x)=\frac{\sum_j\Eta(U_j)P_{U_j}(x_{<\nu})\sum_{x\in A}P_{U_j}(x)}{P_\Eta(x_{<\nu})}\le 1 .
\]
Since $x_\nu$ is independent of $x_{<\nu}$ and $-\log_2 y\ge(1-y)/\ln 2$,
\[
    E_\mu(e_\nu\mid x_{<\nu})
    \;\ge\;\frac{1}{\ln 2}\Bigl(1-\sum_{x:\,P_{U_\mu}(x)>0}P_{U_\mu}(x)\,\frac{q(x)}{P_{U_\mu}(x)}\Bigr)
    \;\ge\;\frac{1}{\ln 2}\Bigl(1-\sum_{x\in A}q(x)\Bigr)\;\ge\;0 .
\]
(This is the non-negativity of the KL divergence from $\mu$ to the sub-probability $q$.) Averaging over $x_{<\nu}$ gives $E_\mu(e_\nu)\ge 0$. For finiteness: $e_\nu$ is bounded below by the integrable function $(1-q(x_\nu)/P_{U_\mu}(x_\nu))/\ln 2$, and by (1) it is bounded above by $-\log_2\Eta(U_\mu)-\sum_{k<\nu}e_k$, which is in turn bounded above by an integrable function using the same lower bound for each $e_k$.

(3) Take expectations in (1) and let $n\to\infty$; the terms are non-negative by (2).

(4) If $m$ of the terms exceed $\varepsilon$ then $m\varepsilon\le\sum_\nu E_\mu(e_\nu)\le-\log_2\Eta(U_\mu)$ by (2) and (3).

(5) The terms of a convergent series tend to zero.
\end{proof}

\begin{remark}[The averaged bound]
Equation \ref{eq:hsi-average-error-bound} is weaker than Equation \ref{eq:hsi-error-bound}, not stronger: dividing by $|X|$ loses the per-term conclusion. (A sequence equal to $1$ on perfect squares and $0$ elsewhere has averages tending to $0$ but does not itself tend to $0$.) Theorem \ref{thm:hsi-rate} is derived from the total bound, and Equation \ref{eq:hsi-average-error-bound} follows from part (1) and Lemma \ref{lem:hsi-chain-rule} by dividing by $|X|$. I suggest demoting it to a remark.
\end{remark}

\begin{remark}[Normalisation]
% NOT formalised in Lean.
$P_{\Eta(\cdot\mid X)}$ is in general a strict sub-probability on $A$. Renormalising it to sum to one over $A$ only raises each $P_{\Eta(\cdot\mid X)}(x)$ and so lowers each $e_\nu$, while $E_\mu(e_\nu)$ remains a KL divergence and so remains non-negative; parts (2)--(5) therefore also hold for the renormalised predictor.
\end{remark}

\begin{remark}[Assumption \ref{ass:hsi-rate}(ii) cannot be dropped]
% Lean: HsiLean/RateCounterexample.lean, theorem rate_counterexample (UNAPPROVED), with weights
% Eta(U_mu) = Eta(U_1) = 1/2. There P_{U_mu} IS the sampling distribution on the items, so the
% negative excess is not caused by a wrong reference: it is caused by the predictor's
% probabilities summing to more than one over A (they tend to 2).
Let $A=\{0,00\}$, which is not prefix-free, and let $\mu$ sample $0$ and $00$ with probability $\tfrac12$ each, so $P_{U_\mu}(0)=P_{U_\mu}(00)=\tfrac12$. Let $U_1$ compute the measure concentrated on $000\dots$, so $P_{U_1}(0)=P_{U_1}(00)=1$, and let $\Eta(U_\mu),\Eta(U_1)>0$. After $n$ samples $P_{U_\mu}(X)=2^{-n}$ and $P_{U_1}(X)=1$, so
\[
    P_{\Eta(\cdot\mid X)}(x)=\frac{\Eta(U_\mu)2^{-(n+1)}+\Eta(U_1)}{\Eta(U_\mu)2^{-n}+\Eta(U_1)}\longrightarrow 1,
    \qquad e_{n+1}\longrightarrow-\log_2\frac{1}{1/2}=-1 .
\]
The excess error converges to $-1$, not $0$: Equation \ref{eq:hsi-error-bound} still holds, but Equation \ref{eq:hsi-convergence} fails, because the ``true'' model is not the best predictor of items of $A$. The paper should therefore say what kind of object a dataset item is (sequences of a fixed length, or any prefix-free set, suffice).
\end{remark}

% ---------------------------------------------------------------------------------------------
% Added 2026-11-11 (Claude, Fable 5.2). UNAPPROVED, unverified by Nathan.
% Lean counterpart: HsiLean/HsiRate.lean (all statements UNAPPROVED; no sorry; axioms
% propext, Classical.choice, Quot.sound only).
\paragraph{The end-to-end result.}
\begin{lemma}[Prefix-free sets]
\label{lem:prefix-free-sum}
% Lean: tsum_le_one_of_prefixFree (finite, relative version: sum_le_of_prefixFree)
Let $A\subseteq\mathbb{B}^*$ be prefix-free, that is, no element of $A$ is a proper prefix of another. Then $\sum_{x\in A}\nu(x)\le\nu(\epsilon)\le 1$ for every semimeasure $\nu$.
\end{lemma}
\begin{proof}
It suffices to treat finite $A$. We show by induction on $n$ that if every element of $A$ extends $y$ by at most $n$ symbols then $\sum_{x\in A}\nu(x)\le\nu(y)$. If $y\in A$ then $A=\{y\}$ since $A$ is prefix-free (this also covers $n=0$). Otherwise every element of $A$ extends $y0$ or $y1$ by at most $n-1$ symbols, so by the induction hypothesis the sum is at most $\nu(y0)+\nu(y1)\le\nu(y)$. Take $y=\epsilon$.
\end{proof}
Lemma \ref{lem:prefix-free-sum} gives Assumption \ref{ass:hsi-rate}(ii) for free, so the two halves of the paper combine:

\begin{theorem}[HSI end to end]
\label{thm:hsi-end-to-end}
% Lean: hsi_end_to_end (general enumeration nu), hsi_end_to_end_enum (the enumeration of all
% enumerable semimeasures, with the truth given as a semimeasure rather than an index).
Let $\{\nu_j\}$ be an effective enumeration of enumerable semimeasures in which every universal mixture appears, and let $\Eta$ be lower-semicomputable with $\Eta(j)>0$ for all $j$ and $\sum_j\Eta(j)\le 1$. Let $A\subseteq\mathbb{B}^*$ be prefix-free, and let $x_1,x_2,\dots$ be sampled independently from some $\nu_\mu$ in the enumeration that is a probability distribution on $A$. Then:
\begin{enumerate}
    \item $P_\Eta=\sum_j\Eta(j)\,\nu_j$ is a universal mixture (Theorem \ref{thm:enums-of-ums=ums});
    \item the conditioned predictor $P_{\Eta(\cdot\mid x_{<\nu})}$ satisfies every part of Theorem \ref{thm:hsi-rate} with constant $-\log_2\Eta(\mu)$; in particular $\sum_{\nu}E_\mu(e_\nu)\le-\log_2\Eta(\mu)$, at most $-\log_2\Eta(\mu)/\varepsilon$ positions have $E_\mu(e_\nu)>\varepsilon$, and $E_\mu(e_\nu)\to 0$.
\end{enumerate}
If $\{\nu_j\}$ enumerates all enumerable semimeasures, this holds for every enumerable $\mu$ that is a probability distribution on $A$.
\end{theorem}
% Differences from the paper's wording: items come from a prefix-free set A; the dataset is an
% ordered i.i.d. sequence (not a set); the weights are over the indices j of the enumeration
% (if mu appears at several indices, the constant uses the weight of one such index, which is
% weaker than using their sum); the expectation is over x_{<nu} and x_nu jointly.
